\documentclass[12pt,a4paper]{article}
%% ═══════════════════════════════════════════════════════════════════
%% VSEPR-SIM Revision Markup — permanent colour definitions
%% DO NOT EDIT these definitions.  They are hardcoded project-wide.
%% ═══════════════════════════════════════════════════════════════════
\usepackage{xcolor}
\usepackage{soul}
\usepackage[normalem]{ulem}
\usepackage{tcolorbox}
\tcbuselibrary{skins,breakable}

%% ── Hardcoded colours ──────────────────────────────────────────────
\definecolor{vsgreen}{HTML}{27AE60}     %% additions
\definecolor{vsred}{HTML}{E74C3C}       %% deletions
\definecolor{vsyellow}{HTML}{F1C40F}    %% edits
\definecolor{vsyellowbg}{HTML}{FEF9E7}  %% edit background
\definecolor{vscyan}{HTML}{1ABC9C}      %% figures
\definecolor{vsdarkred}{HTML}{C0392B}   %% charts

%% ── Inline markup commands ────────────────────────────────────────
\newcommand{\vsadd}[1]{\textcolor{vsgreen}{\textbf{+}~#1}}
\newcommand{\vsdel}[1]{\textcolor{vsred}{\sout{#1}}}
\newcommand{\vsedit}[1]{\sethlcolor{vsyellowbg}\hl{\textcolor{black}{#1}}}
\newcommand{\vsfig}[1]{\textcolor{vscyan}{\textit{[Fig]~#1}}}
\newcommand{\vschart}[1]{\textcolor{vsdarkred}{\textit{[Chart]~#1}}}

%% ── Margin annotations ────────────────────────────────────────────
\newcommand{\vsaddnote}[1]{\marginpar{\tiny\textcolor{vsgreen}{+#1}}}
\newcommand{\vsdelnote}[1]{\marginpar{\tiny\textcolor{vsred}{-#1}}}
\newcommand{\vseditnote}[1]{\marginpar{\tiny\textcolor{vsyellow}{\textrm{#1}}}}

%% ── Block-level environments (tcolorbox) ──────────────────────────
\newtcolorbox{vsaddblock}{%
  colback=vsgreen!8, colframe=vsgreen, title=\textbf{Addition},
  fonttitle=\small, breakable, left=4pt, right=4pt, top=2pt, bottom=2pt}
\newtcolorbox{vsdelblock}{%
  colback=vsred!8, colframe=vsred, title=\textbf{Deletion},
  fonttitle=\small, breakable, left=4pt, right=4pt, top=2pt, bottom=2pt}
\newtcolorbox{vseditblock}{%
  colback=vsyellowbg, colframe=vsyellow, title=\textbf{Edit},
  fonttitle=\small, breakable, left=4pt, right=4pt, top=2pt, bottom=2pt}
\newtcolorbox{vsfigblock}{%
  colback=vscyan!8, colframe=vscyan, title=\textbf{Figure},
  fonttitle=\small, breakable, left=4pt, right=4pt, top=2pt, bottom=2pt}
\newtcolorbox{vschartblock}{%
  colback=vsdarkred!8, colframe=vsdarkred, title=\textbf{Chart},
  fonttitle=\small, breakable, left=4pt, right=4pt, top=2pt, bottom=2pt}
%% ═══════════════════════════════════════════════════════════════════
%% END VSEPR-SIM Revision Markup
%% ═══════════════════════════════════════════════════════════════════

\usepackage{geometry}
\geometry{margin=2.5cm, marginparwidth=2cm}
\usepackage{parskip}
\usepackage{booktabs}

\title{\textbf{VSEPR-SIM Revision Markup} \\ \large Helper Demo}
\author{VSEPR-SIM 3.0.0}
\date{}

\begin{document}
\maketitle

% ─────────────────────────────────────────────────────────────────────────────
\section*{Colour Reference}

\begin{center}
\begin{tabular}{llll}
\toprule
\textbf{Command}   & \textbf{Colour}  & \textbf{Hex}    & \textbf{Meaning}       \\
\midrule
\verb|\vsadd|      & \textcolor{vsgreen}{\textbf{Green}}      & \texttt{27AE60} & Addition / new text    \\
\verb|\vsdel|      & \textcolor{vsred}{\textbf{Red}}          & \texttt{E74C3C} & Deletion / removed     \\
\verb|\vsedit|     & \sethlcolor{vsyellowbg}\hl{Yellow}       & \texttt{F1C40F} & Edit / modified        \\
\verb|\vsfig|      & \textcolor{vscyan}{\textbf{Cyan}}        & \texttt{1ABC9C} & Figure reference       \\
\verb|\vschart|    & \textcolor{vsdarkred}{\textbf{Dark Red}} & \texttt{C0392B} & Chart reference        \\
\bottomrule
\end{tabular}
\end{center}

% ─────────────────────────────────────────────────────────────────────────────
\section*{Inline Markup}

The activation energy of Mo dissolution in FLiBe was \vsdel{220 kJ/mol}
\vsadd{238 kJ/mol} after re-fitting to the updated ORNL-4254 dataset.
\vsaddnote{ref updated}

The Larson--Miller constant was \vsedit{revised from $C=19$ to $C=20$}
based on Hastelloy-N creep data from DeVan \& Evans (1962).
\vseditnote{LM const}

See \vsfig{Figure~3 — Cp(T) curves for INOR-8 alloys}
and \vschart{Chart~2 — stochastic Cp spread at 1000 K} for
the quantitative comparison.
\vsdelnote{fig moved}

% ─────────────────────────────────────────────────────────────────────────────
\section*{Block-level Environments}

\begin{vsaddblock}
Added in revision 2: the cylindrical heat-conduction analysis now includes
a wall heat-flux term $q'' = k(T)\,\nabla T$ evaluated at $r = r_o$.
The result is tabulated in Section~9.
\end{vsaddblock}

\begin{vsdelblock}
Removed: the earlier rule-of-mixtures estimate of thermal conductivity
using a simple arithmetic mean.  This underestimated scattering effects
in W-doped alloys by up to 18\%.
\end{vsdelblock}

\begin{vseditblock}
Revised: the stochastic Monte Carlo sample count was increased from
$n = 100$ to $n = 200$ to reduce the 95\,\% confidence interval on
$\overline{C}_p$ from $\pm 1.4$ to $\pm 0.9$\,J\,mol$^{-1}$\,K$^{-1}$.
\end{vseditblock}

\begin{vsfigblock}
Figure 5 — Radial temperature profile $T(r)$ for a Hastelloy-N tube
($r_i = 10$\,mm, $r_o = 14$\,mm) at steady state under 5\,MW\,m$^{-2}$
wall loading.  Generated by \texttt{cylinder\_steady\_state()}.
\end{vsfigblock}

\begin{vschartblock}
Chart 3 — Stochastic spread of $C_p(T)$ for nominal INOR-8 and
Mo-22\,wt.\% variant.  200 Monte Carlo samples, $3\sigma$ composition
perturbation.  Shaded band = 95\,\% confidence interval.
\end{vschartblock}

% ─────────────────────────────────────────────────────────────────────────────
\section*{Diff-line Pattern}

The following illustrates the diff-line convention used by
\texttt{markup\_diff\_line()} in \texttt{pykernel/latex\_markup.py}
and \texttt{inject\_latex\_markup.sh}:

\begin{verbatim}
+ New sentence added in this revision.
- Old sentence that was removed.
~ Sentence that was modified (shown highlighted).
[Fig] Figure 3 caption was updated.
[Chart] Chart 2 axis labels corrected.
\end{verbatim}

Rendered output:

\vsadd{New sentence added in this revision.}

\vsdel{Old sentence that was removed.}

\vsedit{Sentence that was modified (shown highlighted).}

\vsfig{Figure~3 caption was updated.}

\vschart{Chart~2 axis labels corrected.}

% ─────────────────────────────────────────────────────────────────────────────
\section*{Integration Notes}

\begin{itemize}
  \item \textbf{C/C++ generators} --- \texttt{\#include "include/chart/latex\_markup.h"},
        call \texttt{latex\_markup\_add(buf, sz, text)} etc.,
        then \texttt{latex\_markup\_write\_preamble(fp)} before the document body.
  \item \textbf{Shell batch injection} ---
        \texttt{bash scripts/inject\_latex\_markup.sh file.tex} (in-place) or
        \texttt{--all docs/} for a whole directory.
  \item \textbf{Python report generators} ---
        \texttt{from pykernel.latex\_markup import vsadd, vsdel, vsedit, vsfig, vschart}
        and \texttt{inject\_markup\_preamble("report.tex")}.
  \item \textbf{Colours are permanent} --- hex values are hardcoded in all three
        authorities (C header, Shell script, Python module).  Do not change them.
\end{itemize}

\end{document}
